\section{Arithmetic target and the specific continuation}
The preceding relative-boundary construction forced an available high packet to
be a subgroup norm after reduction modulo two. The latest logarithmic tranche
then improved one sufficient high-factor capacity, but did not force that
inventory to occur. Here we keep the normalized \emph{unit multiplying the
relative norm}. It can change the occupied quotient classes. The result is an
integer lower bound with complete original-word return, and an explicit infinite
family separated from the preceding declared one-/two-prime affine tests.
This is not a new generic theorem about group algebras or a universal ES proof.
The classical weight-retention antecedent is attributed below.

Fix
\[
 k\ge3,\quad u=2^{2k-5},\quad p\equiv1\pmod8,\quad p>2u,
 \qquad N=p+4u=\prod_q q^{e_q}.                            \tag{A1}
\]
The structural results allow integer $p$; prime examples carry separate proofs.
The original fixed-seed middle states correspond bijectively to
\[
 Q\mid N,\quad Q\equiv-1\pmod{2^k},\qquad R=N/Q,\quad a=(p+R)/4.
                                                               \tag{A2}
\]
Indeed $K(u)=2^{k-2}$ and $u\mid a^2\iff K(u)\mid a$.
Because $N\equiv p\pmod{2^k}$, (A2) gives $R\equiv-p\pmod{2^k}$.
Both $Q,R$ are at least three and $3\pmod4$; $N<3p$ gives $Q,R<p$.
Moreover $4(a+u)=R+N$ is divisible by $R$. The converse reverses these steps.
The complete ordered return is
\[
 g=\gcd(a,u),\quad(h,r,s)=(g^2/u,u/g,a/g),\quad
 \lambda=(r+s)/R,\quad(x,y,z)=(a,phs\lambda,phr\lambda),\qquad
 u=\frac{pa^2}{Ry-pa}.                                    \tag{A3}
\]
Valuations prove $a=hrs,u=hr^2,\gcd(r,s)=1$; all are units modulo $R$.
The middle gate gives $R\mid r+s$. Clearing the reciprocals gives
$p\lambda+r+s=(p+R)\lambda=4hrs\lambda$.
These are inherited Type II coordinates \cite[\S2, (2.13)--(2.22)]{ET}; in that
source $(a,b,c,d,e,f)=(r,s,\lambda,h,R,Q)$ and its order is $(x,z,y)$.

Retain a physical packet on disjoint original prime-index blocks:
\[
 v_q(W)=a_q+\sigma_qt_i\ (q\in B_i),\quad0\le t_i\le E_i,
 \qquad W(\mathbf t)=A\prod_iG_i^{t_i}.                    \tag{A4}
\]
Each step vector is nonzero and both endpoints lie in $[0,e_q]$. Thus every
word is an available positive integer divisor, even when $G_i$ is a ratio.
The inverse reads $(v_q(W)-a_q)/\sigma_q$ at a nonzero step and checks the
other coordinates. Use either cyclic chart $H_b=\langle b\rangle$, $b=3,-3$,
of order $o=2^{k-2}$; they comprise residues $1,3\pmod8$ and $1,5\pmod8$.
Assume $A\notin H_b$, $G_i\in H_b$, and keep their logarithms $a_i\in[0,o)$.
The optional role coordinate $\varepsilon\in\{0,1\}$ has multiplier $p^{-1}$;
it is not a prime factor. A packet target is
\[
 Wp^{-\varepsilon}\equiv-1\pmod{2^k}.
 \quad\varepsilon=0:\ Q=W;
 \quad\varepsilon=1:\ R=W,\ Q=N/W.                        \tag{A5}
\]
Fixed-role source words are injective. Forgetting the role has fibres of size
at most two, represented by the complementary words $Q$ and $R$.

\section{The exact unit-corrected relative boundary}
Let $o=2^n$ and $D\mid o$, with $1<D<o$. Put
\[
 A_o=\mathbb F_2[X]/(X^o-1)=\mathbb F_2[T]/T^o,\quad T=X+1,
 \qquad \mathcal N_D=\sum_{j=0}^{o/D-1}X^{jD}=T^{o-D}.
                                                               \tag{B1}
\]
The coordinate change and inverse are both the binomial substitution $X=T+1$.
Let $P\in\mathbb Z_{\ge0}[C_o]$ be the actual integer counting polynomial of
an available labelled packet. Suppose
\[
 0\ne\overline P=T^\nu U(T),\qquad U(0)=1,
 \qquad o-D\le\nu<o.
\]
Define the retained quotient polynomial, in both coordinates, by
\[
 z=\nu-o+D,\quad q_T=T^zU\pmod{T^D},\qquad q(X)=q_T(X+1)\in A_D.
                                                               \tag{B2}
\]
Let $\widehat q$ denote its coefficientwise $0/1$ lift, not a lift of the
$T$-coefficients followed by an unrecorded cancellation.

\begin{theorem}[Integer lower bound and complete boundary image]\label{thm:unit}
There is a unique nonnegative integer polynomial $E$ with
\[
 \boxed{P=\mathcal N_D\widehat q+2E.}                       \tag{B3}
\]
For any nonnegative original complementary packet $G$, define $\pi_DG$ by
summing its coefficients in each congruence class modulo $D$. Then
\[
 \boxed{[X^t](PG)
 =[X^{t\bmod D}](\widehat q\,\pi_DG)+2[X^t](EG)
 \ge[X^{t\bmod D}](\widehat q\,\pi_DG).}                  \tag{B4}
\]
The products on the right use the indicated cyclic rings. A positive right
side supplies an actual original packet word, with an explicit finite inverse.

If $z=0$, multiplication by $q$ is invertible on $A_D$, and
\[
 \iota_q:A_D\xrightarrow{\sim}\ker(T^D:A_o\to A_o),\qquad
 h\longmapsto\mathcal N_Dqh                              \tag{B5}
\]
is a module isomorphism onto its image. Its inverse tests that coefficients
are constant on the $D$-cosets, reads one from each coset, and multiplies by
$q^{-1}$. The ordinary quotient is not this inverse:
$\pi_D\iota_q=(o/D)q=0$.

If $z>0$, the same map has kernel $T^{D-z}A_D$, image $T^\nu A_o$, and
rank $D-z=o-\nu$. A coefficient vector is in that image precisely when it is
constant on $D$-cosets and its readout is divisible by $T^z$.
\end{theorem}
\begin{proof}
Multiplication by $\mathcal N_D$ repeats a coefficient at every member of its
$D$-coset, so $\overline P=\mathcal N_Dq$. Each coefficient of $P$ is at least
its own parity and differs from it by an even nonnegative integer, proving
(B3), including uniqueness. Convolution commutes with this repetition via
$\mathcal N_DF=\iota_D\pi_DF$ over $\mathbb Z$. This proves (B4) with all integer
multiplicities, even when the lower coefficient itself is even.

In $T$-coordinates multiplication by $T^{o-D}$ identifies $A_D$ with
$\ker T^D$. A polynomial is a unit in $A_D$ exactly when its constant
$T$-coefficient is one. Its inverse is
$\sum_{j=0}^{D-1}(1+q_T)^j$. These facts prove (B5) and its stated inverse.
For $z>0$, divide the readout by $T^z$, multiply by $U^{-1}$ modulo $T^{D-z}$,
and retain all additions in $T^{D-z}A_D$. This proves the kernel, image and
all inverse fibres. Summing repeated coefficients gives the ordinary quotient
formula, in which the even integer $o/D$ becomes zero only after reduction.
\end{proof}

This is an actual instance of the SplitZero comparison BR1 \cite{SZ}. For $z=0$
use the prequotient $M=\mathbb F_2[[T]]$, $v=T^D$, and
$f=T^{o-D}U(T)$. The power-series unit is invertible, so $vf(M)=T^oM$.
Both required multipliers are injective \emph{before} quotienting. Using
$\mathbb F_2[T]$ without localization would not justify identifying the ideal
$(T^oU)$ with $(T^o)$. The source and receiving actions of $X=1+T$ commute with
$f$, giving the actual BR2 square. No analytic theta or Gamma norm is imported.

The original carrier remains $G(A)=\{\tau\}\sqcup A^\bullet$, with
$e=0_A^\bullet\ne\tau$. The integer coefficients, presence masks and original
word fibres remain before parity reduction. In particular a positive even
coefficient becomes a supported zero, not an absent word. If $\overline P=0$,
Theorem~\ref{thm:unit} supplies no nonzero layer; it does not assert $P=0$.
Nor is the coefficient-ring inverse a global sparse semiring inverse.

\section{How the actual factor packet and its inverse are computed}
Choose $0\le c_i\le E_i$ for nonidentity cyclic steps, and reserve the complete
coordinate for that choice. Write $w_i=2^{v_2(a_i)}$ and suppose
\[
 \nu=\sum_i c_iw_i<o.                                    \tag{C1}
\]
The first nonzero term of $1+(1+T)^{a_i}$ has degree $w_i$ and coefficient
one. Thus
\[
 U(T)=\prod_i\left(\frac{1+(1+T)^{a_i}}{T^{w_i}}\right)^{c_i},
 \qquad \overline P=T^\nu U(T).                           \tag{C2}
\]
For $c=\sum c_j2^j$, let $t\preceq c$ mean binary-submask containment.
The \emph{integer} packet is
\[
 P(X)=\prod_i\left(\sum_{t_i\preceq c_i}X^{a_it_i}\right). \tag{C3}
\]
Every allowed original parameter occurs once. The equality with the binomial
product is only asserted after reduction modulo two. Repeated copies of one
prime have not been turned into distinguishable coloured primes.

An untouched coordinate ($c_i=0$) can supply its full original interval to
$G$. A partially used coordinate is not also treated as a new independent
remaining factor. More general decompositions require their complete fibre
weights and are not presumed here. For an actual affine packet, (B4) bounds
marked role words; forgetting the role divides the guaranteed count by at
most two, rounded upwards.

To reconstruct, choose an actual $G$-word $g$ contributing to a positive lower
coefficient, with $q_{t-g}=1$. The high coefficient of $P$ at that \emph{fine}
target is odd. Expand (C3) into binary factors and form suffix products over
$\mathbb F_2[C_o]$. At each step the two possible suffix coefficients have
sum one, so exactly one is odd. Select it and continue. The final exponent is
zero; the selected bits return $t_i\preceq c_i$. Use (A4)--(A5) and (A3).
Conversely, retain the affine pattern, role, selected capacities and $G$-word;
valuation inversion and submask membership recover precisely its source word.
All alternative lower words and the complete parity preimages remain fibres,
not an unspecified inverse kernel.

For each finite original factorization, searching capacities, proper indices
and available complementary lower words terminates. The implemented canonical
search uses the same disjoint singleton/same-coset-pair affine patterns as its
predecessor. It is complete for its explicit enumeration, not for arbitrary
larger blocks. Computing $U\bmod T^D$ need not expand the high polynomial, but
the demonstrated suffix inverse may use the full $o$-coordinate ring. No
polynomial-time factorization or subgroup-size-free inverse is claimed.

\paragraph{Classical support bound, proved here in the needed binary form.}
For nonzero $F\in A_{2^n}$ of exact $T$-order $\nu$,
\[
 \operatorname{wt}_X(F)\ge2^{\operatorname{popcount}(\nu)}.\tag{C4}
\]
Write $F=A+X^mB=(A+B)+T^mB$, $m=2^{n-1}$, with $A,B$ of degree below $m$.
If $\nu<m$, apply induction to $A+B$ and
$\operatorname{wt}(A)+\operatorname{wt}(B)\ge\operatorname{wt}(A+B)$.
If $\nu\ge m$, then $A=B$; the weight doubles and induction applies to $A$.
This is the binary weight-retention property of Massey--Costello--Justesen
\cite{MCJ}, not a novelty claim. Their primary repository abstract was read;
the proof above does not depend on an unread detailed theorem numbering.
The bound counts occupied cosets; the target location still requires the
actual $q$, rather than only its weight.

\section{An integer inverse and its nontrivial arithmetic index}
The following explicit family measures what parity alone forgets. Let $D$ be
a power of two and $h$ an odd integer, $1\le h<D$. Write
\[
 q=X^rS_h(X),\quad S_h=1+X+\cdots+X^{h-1},\quad
 a=h^{-1}\bmod D\ (1\le a<D),\quad b=(ha-1)/D.
\]
Write $J_D=1+X+\cdots+X^{D-1}$ for the full norm in $\mathbb Z[C_D]$,
distinct from the relative norm $\mathcal N_D$ in (B1). Geometric multiplication gives
\[
 S_h(X)S_a(X^h)=1+bJ_D\quad\text{in }\mathbb Z[C_D].
\]
Therefore the complete rational inverse is
\[
 \boxed{q^{-1}=X^{-r}\left(S_a(X^h)-\frac bh J_D\right).} \tag{D1}
\]
The image of its integer circulant $C_q$ is exactly
\[
 \boxed{C_q\mathbb Z^D
 =\{y\in\mathbb Z^D:\textstyle\sum y_i\equiv0\pmod h\}.} \tag{D2}
\]
Necessity follows from the column sum $h$. For sufficiency, (D1) is integral
on precisely that stated sufficient condition. Hence the cokernel is
$\mathbb Z/h$, and $|\det C_q|=h$. The source coefficient norm transports to
\[
 \boxed{(\mathcal N_D C_q)^*(\mathcal N_D C_q)
 =(o/D)C_q^{\mathsf T}C_q,}                               \tag{D3}
\]
when the left norm repeats into a larger cyclic group of order $o$.
No normalization changes the original lattice index or cross terms.
For $D=8$, $q=X^3+\cdots+X^7$, and $o=16$, this declared coefficient Gram's first row is
$(10,8,6,4,4,4,6,8)$ and the index is five. Modulo two the map is invertible;
over the integer lattice it is not unimodular. Its rational inverse has
negative coefficients and is not used as a positivity-preserving map.

\paragraph{The original word metric and its actual inverse domain.}
Equation (D3) uses the declared cyclic \emph{coefficient} norm. It is not
silently the quotient norm of the original word space. If the latter has one
orthonormal coordinate per actual word and fibre size $P_r$ at residue $r$,
its residue-sum quotient has norm $\sum_{P_r>0}|y_r|^2/P_r$. Its canonical
section assigns $y_r/P_r$ to every word in that fibre. The kernel is the
fibre-sum-zero space, whose complete star primitive is the vector of
nonreference coefficients; subtracting the canonical section proves the
Pythagorean identity.

A formal unit inverse also does not create words at absent residues. The exact
available domain of the periodic comparison over $\mathbb Q$ is
\[
 \mathcal D_P=\{h:(C_qh)_{r\bmod D}=0\text{ whenever }P_r=0\}.
                                                               \tag{D4}
\]
On this explicitly specified subspace the lift to original words is
$(C_qh)_{r\bmod D}/P_r$ at each word of fibre $r$, and its Gram is
\[
 \boxed{C_q^{\mathsf T}\operatorname{diag}\!\left(
 \sum_{r\equiv j\ (D),\,P_r>0}\frac1{P_r}\right)_{j=0}^{D-1} C_q.} \tag{D5}
\]
The proof is the sum of the squared displayed coefficients, with all original
multiplicities. For (D3)'s $D=8$ example below, $P_8=P_9=P_{10}=0$; (D4) is
exactly $(C_qh)_0=(C_qh)_1=(C_qh)_2=0$, of dimension five. The unrestricted
formal module has dimension eight. The successful vector $h=e_0$ is in the
available domain, but arbitrary translates need not be. This domain restriction
is a calculated linear image, not a universal occupancy assumption.

\section{An unbounded strict arithmetic family}
Let $k\ge6$, $d=2^{k-4}$, $o=4d$, and $u=2^{2k-5}$. Choose distinct actual
primes $q,r,s$, all different from three, in the reduced classes
\[
 q\equiv3^{3d-1},\qquad r\equiv-3^{5d/2},\qquad
 s\equiv-3^{d/2+2}\pmod{2^k}.                              \tag{E1}
\]
Set
\[
 N=3^{d-1}qrs,\qquad p=N-4u,
\]
and choose the primes large enough that $p>2u$. Then
$p\equiv3^{3d}\pmod{2^k}$ and $p\equiv1\pmod{24}$.
The role multiplier $p^{-1}$ has logarithm $d$ in the $3$-chart.
The actual packet with constant $s$ is
\[
 W=3^iq^js,
 \quad0\le i\le d-1,\quad j\in\{0,1\},\quad\varepsilon\in\{0,1\}.
\]
Its generating polynomial, with both roles retained, is
\[
 P_d=(1+\cdots+X^{d-1})(1+X^{3d-1})(1+X^d).
\]
Direct coefficient comparison gives the integer identity
\[
 \boxed{
 P_d=\mathcal N_{2d}\underbrace{(X^{d-1}+\cdots+X^{2d-1})}_{Q_d}
       +2(1+X+\cdots+X^{d-2}).}                           \tag{E2}
\]
The quotient $Q_d$ is a non-monomial unit with $d+1$ occupied coefficients;
its integer circulant has index $d+1$ by (D2).
The fine target for $s$ is $7d/2-2$, whose residue $3d/2-2\pmod{2d}$ occurs
in $Q_d$. Theorem~\ref{thm:unit} therefore forces an original word.

\begin{theorem}[Unique original state and strict unbounded word length]
For every integer parameter in (E1), the complete fixed-seed branch has exactly
one cofactor:
\[
 \boxed{Q=3^{d/2-1}qs,\qquad R=3^{d/2}r.}                 \tag{E3}
\]
Both roles have minimum prime-occurrence length
\[
 \boxed{\Omega(Q)=\Omega(R)=d/2+1=2^{k-5}+1.}              \tag{E4}
\]
The preceding canonical full-chart and stride-relative tests, using disjoint
singleton and same-coset two-prime lines, fail on the entire family, in both
charts. This does not include arbitrary larger-arity affine packets.
\end{theorem}
\begin{proof}
Only $r,s$ are outside the $3$-chart, and each occurs once. A cofactor must
therefore contain exactly one. The inside word logs occupy
$[0,d-1]\cup[3d-1,4d-2]$. The $r$ target is $3d/2$, in neither interval.
The $s$ target is $7d/2-2$, occurring only with $j=1,i=d/2-1$.
This proves the complete cofactor census and (E3)--(E4).

Here is the uniform check against the precise preceding test class. Let its
common lower stride be $D_0=2^v$, with high group order $o/D_0$. For a step
log $a$ and capacity $E$, its maximum high contribution is
\[
 \left\lfloor\frac{E}{\delta}\right\rfloor
 \operatorname{lowbit}\left((a\delta/D_0)\bmod(o/D_0)\right),
 \qquad\delta=D_0/\gcd(a,D_0).
\]
Lower digits can only reduce this capacity. Over every old pattern in either
chart, the total is at most
\[
 \begin{array}{c|cccc}
 D_0&1&2&4\le D_0\le d&2d\\\hline
 \text{maximum upper bound}&2d+2&d+1&2d/D_0&0.
 \end{array}                                               \tag{E5}
\]
Each entry is strictly below the required total $4d/D_0-1$.

For completeness, the finite pattern analysis proving (E5) is independent of
$d$. In the $3$-chart the two inside singleton weights/capacities are
$(1,d-1),(1,1)$; outside singletons have zero capacity. Their inside pair has
weight $d$ (parallel) or $2$ (exchange), capacity one. The outside pair has
weight $2$ when it has nonzero capacity, again at most one. The role has
weight $d$, capacity one. If $D_0\ge4$, all capacity-one weight-one or
weight-two items disappear; only the role, a possible weight-$d$ pair, or
at most $d/D_0-1$ singleton units remain. At strides one and two the same
list gives $2d+2$ and $d+1$.

In the $-3$-chart every prime is outside. The only possible nonzero singleton
capacity is $d/2-1$ at weight two, from the squared $3$ step. The $(3,q)$ pair
has weight $d$ or two, each other pair involving one of these and $r$ or $s$
has odd weight-one log, and the $(r,s)$ pair has weight two. Every pair has
capacity at most one, and disjoint blocks use each prime index at most once.
Together with the same role these give exactly the bounds (E5).
Thus the old necessary total-degree condition fails before any quotient target
can be checked. This proves failure of the declared old test, not nonexistence
of all possible affine maps.
\end{proof}

There is a surviving larger-arity correspondence. The \emph{three-prime} shear
\[
 (v_3,v_r,v_s)=(a+2t,t,1-t),\quad0\le a\le d-3,\quad t=0,1, \tag{E6}
\]
has available endpoints, inverse $t=v_r$, $a=v_3-2v_r$, and step
$3^2r/s\equiv3^{2d}$, a half-turn. Retain $q$ with exponent one and choose
$a=d/2-1$. Its $t=0$ endpoint is (E3). This map reuses the $3$ coordinate
coherently across its two parameters; it is not a disjoint pair block.
It both checks the positive return and explains precisely the scope of the
separation theorem.

Dirichlet's theorem supplies infinitely many choices of the auxiliary primes
in (E1) \cite[Theorem 18.1]{Dirichlet}. It is not asserted that infinitely many
of the derived $p$ are prime. The integer parameters can all be required to
lie in $1\pmod{840}$: choose $r,s$ avoiding $5,7$, then prescribe
$q\equiv(1+4u)(3^{d-1}rs)^{-1}\pmod{35}$ in addition to (E1).
This is a reduced class because $1+2^{2k-3}$ is nonzero modulo both $5$ and $7$.
Together with the already proved $p\equiv1\pmod{24}$ it gives the claim.

\section{Prime examples, finite comparison, and proof limits}
Three actual prime members have the following complete cofactors:
\[
\begin{array}{c|r|rrr|r|r}
 k&p&q&r&s&Q&R\\\hline
 6&1721521&59&23&47&8319&207\\
 7&1492567108321&17099&239&167&77099391&19359\\
 8&2413105093468609&1259&223&599&1649306367&1463103
\end{array}
\]
The minimum lengths in both roles are three, five and nine. Thus the last
example defeats an eight-occurrence search in this branch, not in every seed
of that prime. Their complete-order primality proofs use bases $22,17,23$ and
respectively the fully factored predecessors
\[
\begin{aligned}
1721520&=2^4 3^3 5\cdot797,\\
1492567108320&=2^5 3\cdot5\cdot281\cdot11065889,\\
2413105093468608&=2^6 3\cdot7\cdot11\cdot61\cdot887\cdot3016691.
\end{aligned}
\]
The checker proves primality of all factor primes by trial division and checks
$a^{p-1}=1\pmod p$ and $\gcd(a^{(p-1)/q}-1,p)=1$ for every factor prime $q$.
Every prime divisor of $p$ must then be $1\pmod{p-1}$, proving $p$ prime.
No probable-prime test is used as the certificate.

For the first example, the full relative quotient is
$X^3+\cdots+X^7$ in $A_8$, not one; its target coefficient is at four.
The normalized original state is
\[
 (p,R,u,h,r,s,\lambda)=(1721521,207,128,8,4,13451,65),
\]
\[
 (x,y,z)=(430432,12041213064920,3580763680).
\]
For the largest example it is
\[
 (u,a,h,r,s,\lambda)=(2048,603276273732928,2,32,9426191777077,6442603),
\]
\[
 (x,y,z)=(603276273732928,
 293091938807452961286828738688138958,
 994987399327753007950528).
\]
All original prime valuations, roles, complete divisor lists and ordered
inverses are retained in the JSON.

\paragraph{The finite comparison is not the universal claim.}
At all 4519 hard-class primes through two million there are 37289 valid
$k\ge4$ seed branches and 4583 original states in 3097 occupied branches.
The preceding canonical relative class certifies 1723 branches on 1595 primes;
its union with the new test certifies 1778 branches on 1627 primes.
The stronger previous union (relative, two-prime words, and $3,11$ capacity)
certifies 3073 branches on 2402 primes; adding the new test gives 3074 branches
on 2403 primes. The strict addition is $(1721521,6)$. A three-occurrence search
combined with the previous tests already covers all occupied branches in this
small range. Therefore no advantage over that stronger finite baseline, or new
overall ES verification range, is claimed. The unbounded family, rather than
the small census alone, justifies retaining the unit correction.

The source-ready theorem is the integer lower identity (B4), with an actual
bounded inverse and an explicit factor family. For an arbitrary prescribed
prime, its required packet inventory and a positive target coefficient are
still to be forced in at least one seed. Keeping $U$ prevents the map from
changing that target; it does not itself prove universal occupancy.
The next arithmetic direction exposed by (E6) is a coordinated higher-arity
factor exchange, not an unconstrained ratio of unavailable factors.
